% =====================================================================
% CAPÍTULO 6: DESARROLLO DEL APLICATIVO WEB
% Versión académica
% =====================================================================

\chapter{Desarrollo del Aplicativo Web}
\label{ch:desarrollo}

% --- Ajustes editoriales locales del capítulo: sangría, espaciado y control de desbordamientos ---
\setlength{\parindent}{0pt}
\setlength{\parskip}{6pt}
\sloppy
\emergencystretch=2em
\raggedbottom
\section{Enfoque funcional del sistema}

El aplicativo web se concibe y construye como una plataforma de grado de producción diseñada para centralizar, controlar y auditar el ciclo de vida documental y operativo de las órdenes de trabajo de CERMONT S.A.S. 

\subsection{Matriz de estado de implementación y alcance}

Para preservar la integridad académica y la transparencia técnica ante el jurado, la Tabla \ref{tab:matriz-estado} detalla el estado actual de las funcionalidades descritas en el sistema v1.0. Se distinguen las capacidades implementadas con pruebas de aquellas propuestas como evoluciones futuras.

\begin{table}[!htbp]
\centering
\caption{Matriz de estado de implementación. Fuente: Elaboración propia.}
\label{tab:matriz-estado}
\scriptsize
\renewcommand{\arraystretch}{0.92}
\def\EstadoRow#1#2#3{\parbox[t]{4.4cm}{\raggedright #1} & \parbox[t]{3.3cm}{\raggedright #2} & \parbox[t]{5.3cm}{\raggedright #3}\\\hline}
\begin{tabular}{|l|l|l|}
\hline\ClassicTableHeader
\textbf{Funcionalidad} & \textbf{Estado} & \textbf{Evidencia / Justificación} \\
\hline
\EstadoRow{Autenticación JWT + RBAC}{\textbf{Implementado}}{Controladores y middlewares en \path{backend/src/controllers/auth.controller.ts}, \path{backend/src/middlewares/auth.middleware.ts} y roles en \path{packages/domain/src/roles.ts}.}
\EstadoRow{FSM multietapa de órdenes}{\textbf{Implementado}}{Estados y transiciones en \path{packages/shared-types/src/schemas/work-order-fsm.ts} y \path{backend/src/services/order/order-rules.ts}.}
\EstadoRow{Operación offline con Service Worker, Dexie e IndexedDB}{\textbf{Parcialmente implementado con evidencia E2E}}{Service Worker fuente en \path{frontend/src/app/sw.ts}, ruta generada en \path{/serwist/sw.js}, persistencia local en \path{frontend/src/lib/offline/offline-db.ts}, cola de sincronización en \path{frontend/src/lib/offline/sync-engine.ts} y endpoints backend en \path{backend/src/modules/sync/}. La prueba de producción verifica app shell y rutas internas calentadas; la captura offline completa de cada formulario sigue requiriendo validación por corte vertical.}
\EstadoRow{Generación PDF}{\textbf{Implementado}}{Servicios \path{backend/src/services/report.service.ts} y \path{backend/src/services/pdf-generator.service.ts}.}
\EstadoRow{Evidencias con georreferenciación opcional}{\textbf{Implementado}}{Contrato en \path{packages/shared-types/src/schemas/evidence.schema.ts} y servicio en \path{backend/src/services/evidence.service.ts}.}
\EstadoRow{Seguimiento de cierre administrativo}{\textbf{Implementado como seguimiento documental interno}}{Registro de actas, SES, facturas y pagos; no integra SAP Ariba ni emite facturación electrónica DIAN.}
\EstadoRow{Dashboard y analítica operativa}{\textbf{Implementado}}{Contrato en \path{packages/shared-types/src/schemas/analytics.schema.ts} y vistas de panel en \path{frontend/src/app/(dashboard)/dashboard/}.}
\EstadoRow{Asistente AI (Chat)}{\textbf{Complementario / no crítico}}{Módulo \path{ai/}; no se utiliza como evidencia principal de cumplimiento de objetivos.}
\EstadoRow{OCR / Extracción de Datos}{\textbf{Propuesto v2.0}}{Requiere integración con API de visión.}
\EstadoRow{Firma Digital Criptográfica}{\textbf{Propuesto v2.0}}{Requiere certificado digital (PKI).}
\EstadoRow{Integración SAP Ariba}{\textbf{Fuera de alcance}}{Requiere contrato SAP Enterprise API.}
\EstadoRow{Emisión de facturación electrónica DIAN}{\textbf{Fuera de alcance}}{El sistema registra seguimiento documental; la emisión requiere software contable certificado.}
\hline
\end{tabular}
\end{table}

Esta delimitación asegura que el sistema entregado se enfoque en resolver la fragmentación documental interna, dejando las integraciones con sistemas externos de terceros como trabajos futuros de integración empresarial.

\begin{figure}[!htbp]
\centering
\resizebox{0.95\textwidth}{!}{%
\begin{tikzpicture}[node distance=0.9cm, every node/.style={font=\scriptsize}, box/.style={draw, rounded corners, align=center, minimum width=2.1cm, minimum height=0.75cm, fill=blue!6}]
\node[box] (sol) {Solicitud};
\node[box, right=of sol] (prop) {Propuesta\\/ PO};
\node[box, right=of prop] (plan) {Planeación\\validada};
\node[box, right=of plan] (exe) {Ejecución\\+ evidencias};
\node[box, right=of exe] (rep) {Informe\\y acta};
\node[box, right=of rep] (adm) {SES\\Factura\\Pago};
\draw[-{Latex[length=2mm]}] (sol)--(prop); \draw[-{Latex[length=2mm]}] (prop)--(plan); \draw[-{Latex[length=2mm]}] (plan)--(exe); \draw[-{Latex[length=2mm]}] (exe)--(rep); \draw[-{Latex[length=2mm]}] (rep)--(adm);
\node[draw, rounded corners, fill=green!8, align=center, below=0.7cm of exe, minimum width=8.5cm] {Cada transición exige datos, documentos, permisos y auditoría};
\end{tikzpicture}
%
}
\caption{Orquestación del flujo de trabajo digitalizado (BPMN) que integra los 14 pasos operativos en el aplicativo. Fuente: elaboración propia.}
\label{fig:bpmn-digitalizado}
\end{figure}

\begin{figure}[!htbp]
\centering
\resizebox{0.95\textwidth}{!}{%
\begin{tikzpicture}[
  every node/.style={font=\tiny},
  stepbox/.style={draw=institucionalBlue, rounded corners=3pt, fill=institucionalGrey!35, align=center, minimum width=1.62cm, minimum height=0.62cm, line width=0.65pt},
  header/.style={font=\scriptsize\bfseries, text=institucionalBlue},
  arrow/.style={-{Latex[length=1.6mm]}, draw=institucionalBlue, line width=0.65pt},
  bridge/.style={-{Latex[length=1.6mm]}, draw=institucionalRed, line width=0.85pt}
]
\foreach \x/\txt in {0/Solicitud,1/Visita,2/Propuesta,3/PO,4/Planeación,5/Ejecución,6/Informe}{\node[stepbox] (a\x) at (\x*2,1.2) {\txt};}
\foreach \x/\txt in {0/Acta,1/Firma,2/SES,3/SES aprob.,4/Factura,5/Fact. aprob.,6/Pago}{\node[stepbox] (b\x) at (\x*2,0) {\txt};}
\foreach \x in {0,...,5}{\pgfmathtruncatemacro{\y}{\x+1}\draw[arrow] (a\x)--(a\y);\draw[arrow] (b\x)--(b\y);}
\draw[bridge] (a6)--(b0);
\node[anchor=east,header] at (-0.4,1.2) {Operativo};
\node[anchor=east,header] at (-0.4,0) {Administrativo};
\end{tikzpicture}
%
}
\caption{Flujo operativo-documental implementado para CERMONT S.A.S. Fuente: Elaboración propia.}
\label{fig:flujo-proceso}
\end{figure}

\section{Módulos del sistema: Cobertura del flujo operativo y de soporte}

El aplicativo organiza sus funcionalidades en módulos API agrupados por dominio de negocio. Para evitar ambigüedades entre módulos de negocio y grupos de rutas, el documento distingue entre módulos operativos vinculados a los 14 pasos de CERMONT S.A.S. y módulos transversales que respaldan seguridad, sincronización, auditoría y observabilidad de la plataforma.

\begin{table}[!htbp]
\centering
\caption{Módulos operativos principales (Pasos 1-14). Fuente: Elaboración propia.}
\label{tab:modulos-operativos}
\small
\def\ModuloRow#1#2#3{\parbox[t]{3.4cm}{\raggedright #1} & \parbox[t]{5.0cm}{\raggedright #2} & \parbox[t]{5.0cm}{\raggedright #3}\\\hline}
\begin{tabular}{|l|l|l|}
\hline\ClassicTableHeader
\textbf{Módulo} & \textbf{Función Técnica} & \textbf{Paso CERMONT} \\
\hline
\ModuloRow{Autenticación}{Control de acceso RBAC y gestión de sesión JWT HttpOnly.}{Transversal}
\ModuloRow{Órdenes de trabajo}{Motor FSM central de 15 estados y entidad raíz de trazabilidad.}{1 - 14}
\ModuloRow{Propuestas y PO}{Generación de propuestas económicas y validación de órdenes de compra.}{3 - 4}
\ModuloRow{Planeación y Kits}{Reserva de recursos, herramientas y personal técnico.}{5}
\ModuloRow{Checklists y Campo}{Ejecución de inspecciones y formularios con capacidad offline.}{6}
\ModuloRow{Evidencias}{Captura de soportes fotográficos geo-localizados.}{6}
\ModuloRow{Informes Técnicos}{Generación automatizada de reportes PDF (pdf-lib).}{7}
\ModuloRow{Cierre Consolidado}{Agregación de datos de acta, SES, factura y pago (Read Model).}{8 - 14}
\hline
\end{tabular}
\end{table}

Adicionalmente, el sistema se apoya en módulos de soporte como \texttt{analytics} (dashboard), \texttt{sync} (sincronización offline), \texttt{audit} (logs inmutables), \texttt{users}, \texttt{documents}, \texttt{maintenance}, \texttt{inspections}, \texttt{history}, \texttt{alerts}, \texttt{admin}, \texttt{assets} y \texttt{ai} (asistente complementario).

\section{Evidencia visual del aplicativo desarrollado}
\label{sec:evidencia-visual-aplicativo}

Las siguientes capturas documentan pantallas reales del aplicativo y se incorporan como evidencia visual del desarrollo. Cada pantalla se interpreta desde su función operativa y desde su relación con el flujo de 14 pasos de CERMONT S.A.S. Cuando una captura muestra un estado vacío, este se entiende como evidencia de la interfaz, filtros y controles disponibles, no como medición de uso ni como resultado operativo cuantificado.

\begin{figure}[H]
\centering
\AppScreenshot[width=0.96\textwidth]{imagenes/KPIS.png}
\caption{Panel de control y visualización del flujo operativo de 14 pasos. Fuente: captura propia del aplicativo CERMONT.}
\label{fig:captura-dashboard-14-pasos}
\end{figure}

La Figura \ref{fig:captura-dashboard-14-pasos} muestra el panel de control principal, donde se presenta una lectura general del estado operativo y una representación navegable del ciclo de 14 pasos. Esta pantalla cumple la función de punto de entrada para gerencia y coordinación, ya que permite ubicar el avance de solicitudes, propuestas, planeación, ejecución, informes, actas, SES, facturación y pagos dentro de un mismo entorno. Su aporte al proyecto consiste en evidenciar que el aplicativo no se diseñó como un tablero genérico, sino como una interfaz orientada al proceso real de CERMONT S.A.S.

\begin{figure}[H]
\centering
\AppScreenshot[width=0.96\textwidth]{imagenes/captura_usuarios_roles.png}
\caption{Gestión de usuarios, roles y estado de acceso. Fuente: captura propia del aplicativo CERMONT.}
\label{fig:captura-usuarios-roles}
\end{figure}

La Figura \ref{fig:captura-usuarios-roles} evidencia el módulo de administración de usuarios y roles. Esta funcionalidad permite asociar cada cuenta con un perfil de responsabilidad, controlar su estado y soportar la aplicación del modelo RBAC descrito en la arquitectura. En el flujo operativo, este módulo es transversal, porque determina qué usuarios pueden crear solicitudes, aprobar información, ejecutar actividades, registrar evidencias o intervenir en el cierre administrativo.

\begin{figure}[H]
\centering
\AppScreenshot[width=0.96\textwidth]{imagenes/captura_nueva_solicitud.png}
\caption{Formulario de creación de solicitud de trabajo. Fuente: captura propia del aplicativo CERMONT.}
\label{fig:captura-nueva-solicitud}
\end{figure}

La Figura \ref{fig:captura-nueva-solicitud} corresponde al registro inicial de una solicitud de trabajo. El formulario recoge datos básicos del cliente, ubicación, contacto, prioridad y descripción del requerimiento. Esta pantalla se relaciona con el Paso 1 del proceso CERMONT, porque formaliza la entrada del servicio al sistema y evita que la necesidad del cliente quede dispersa en correos, llamadas o mensajes sin trazabilidad documental.

\begin{figure}[H]
\centering
\AppScreenshot[width=0.96\textwidth]{imagenes/PROPUESTA.png}
\caption{Creación de una propuesta económica. Fuente: captura propia del aplicativo CERMONT.}
\label{fig:captura-nueva-propuesta}
\end{figure}

La Figura \ref{fig:captura-nueva-propuesta} presenta la pantalla de creación de propuestas, con campos para información del cliente, vigencia e ítems económicos. La funcionalidad se vincula con el Paso 3 del flujo, en el cual se estructura la oferta técnica y económica que será evaluada por el cliente. En términos operativos, esta vista aporta evidencia de la captura estructurada de valores, cantidades y descripciones, datos que posteriormente alimentan la trazabilidad de costos y cierre.

\begin{figure}[H]
\centering
\AppScreenshot[width=0.96\textwidth]{imagenes/captura_listado_propuestas.png}
\caption{Listado y seguimiento de propuestas económicas. Fuente: captura propia del aplicativo CERMONT.}
\label{fig:captura-listado-propuestas}
\end{figure}

La Figura \ref{fig:captura-listado-propuestas} muestra el listado de propuestas, sus estados y filtros de consulta. Esta vista permite al área comercial revisar propuestas creadas, aprobadas o pendientes, y se relaciona con la transición entre el Paso 3 y el Paso 4. Su valor como evidencia radica en mostrar que la propuesta no queda como un documento aislado, sino como un registro consultable dentro del flujo que habilita la aprobación con orden de compra.

\begin{figure}[H]
\centering
\AppScreenshot[width=0.96\textwidth]{imagenes/captura_ordenes_trabajo.png}
\caption{Consulta y seguimiento de órdenes de trabajo. Fuente: captura propia del aplicativo CERMONT.}
\label{fig:captura-ordenes-trabajo}
\end{figure}

La Figura \ref{fig:captura-ordenes-trabajo} evidencia el módulo de órdenes de trabajo, entidad central para articular los pasos operativos y administrativos. La pantalla permite consultar órdenes, filtrar registros y revisar su estado. Dentro del flujo de 14 pasos, la orden actúa como hilo conductor entre solicitud, visita, propuesta, planeación, ejecución, evidencias, informe, acta, SES, factura y pago.

\begin{figure}[H]
\centering
\AppScreenshot[width=0.96\textwidth]{imagenes/captura_planeacion.png}
\caption{Paquetes de planeación asociados a órdenes de trabajo. Fuente: captura propia del aplicativo CERMONT.}
\label{fig:captura-planeacion}
\end{figure}

La Figura \ref{fig:captura-planeacion} corresponde al módulo de planeación. En esta pantalla se centraliza la preparación de recursos, responsables y documentos requeridos antes de iniciar la actividad. Se relaciona con el Paso 5 del flujo y atiende una de las fallas críticas del proceso: la planeación incompleta de herramientas, personal, permisos y soportes previos a la ejecución.

\begin{figure}[H]
\centering
\AppScreenshot[width=0.96\textwidth]{imagenes/captura_catalogo_kits.png}
\caption{Catálogo de kits típicos para planeación de recursos. Fuente: captura propia del aplicativo CERMONT.}
\label{fig:captura-catalogo-kits}
\end{figure}

La Figura \ref{fig:captura-catalogo-kits} muestra el catálogo de kits, diseñado para reutilizar configuraciones típicas de herramientas, equipos o recursos por tipo de actividad. Esta funcionalidad complementa el Paso 5, ya que permite que la planeación no dependa exclusivamente de memoria individual o listados informales. La captura aporta evidencia de un mecanismo de estandarización operativo, útil para mitigar omisiones en la preparación de trabajos.

\begin{figure}[H]
\centering
\AppScreenshot[width=0.96\textwidth]{imagenes/EJECUCION.png}
\caption{Sesiones de ejecución en campo y estado de sincronización. Fuente: captura propia del aplicativo CERMONT.}
\label{fig:captura-ejecucion}
\end{figure}

La Figura \ref{fig:captura-ejecucion} presenta el módulo de sesiones de ejecución. La pantalla agrupa estados como activas, listas, completadas y sincronización pendiente, además de opciones para cargar documentos, hojas de cálculo y fotografías. Esta vista se relaciona con el Paso 6 y evidencia el soporte del aplicativo para capturar información durante la actividad en campo, incluyendo escenarios de conectividad variable.

\begin{figure}[H]
\centering
\AppScreenshot[width=0.96\textwidth]{imagenes/EVIDENCIAS.png}
\caption{Gestor visual de evidencias por orden y etapa. Fuente: captura propia del aplicativo CERMONT.}
\label{fig:captura-evidencias}
\end{figure}

La Figura \ref{fig:captura-evidencias} corresponde al gestor de soportes visuales. La pantalla permite filtrar evidencias por orden y etapa, así como distinguir registros antes, durante y después de la intervención. Esta funcionalidad se vincula con los Pasos 6 y 7, porque las fotografías y soportes de campo son insumos para el informe técnico y para el cierre documental. La captura evidencia la intención de asociar cada soporte a una orden, evitando que las fotografías queden dispersas sin contexto.

\begin{figure}[H]
\centering
\AppScreenshot[width=0.96\textwidth]{imagenes/captura_gestion_documentos.png}
\caption{Gestión de documentos y soportes del aplicativo. Fuente: captura propia del aplicativo CERMONT.}
\label{fig:captura-gestion-documentos}
\end{figure}

La Figura \ref{fig:captura-gestion-documentos} muestra el módulo de gestión documental. Esta pantalla organiza documentos por tipo, estado y búsqueda, lo que permite relacionar soportes con órdenes, informes, actas o cierre administrativo. Su relación con el proyecto es transversal: los documentos son la evidencia que conecta la ejecución técnica con la aceptación del cliente y con los trámites de facturación.

\begin{figure}[H]
\centering
\AppScreenshot[width=0.96\textwidth]{imagenes/captura_actas_entrega.png}
\caption{Actas de entrega y seguimiento de aceptación. Fuente: captura propia del aplicativo CERMONT.}
\label{fig:captura-actas-entrega}
\end{figure}

La Figura \ref{fig:captura-actas-entrega} evidencia el módulo de actas de entrega. Esta funcionalidad se relaciona con los Pasos 8 y 9, donde se genera el acta, se envía al cliente y se registra su aceptación o firma. La pantalla aporta evidencia de que el cierre técnico no se limita a completar la ejecución, sino que requiere un soporte documental formal para habilitar los pasos administrativos posteriores.

\begin{figure}[H]
\centering
\AppScreenshot[width=0.96\textwidth]{imagenes/PAGOS.png}
\caption{Seguimiento de SES y soportes de Ariba. Fuente: captura propia del aplicativo CERMONT.}
\label{fig:captura-ses-ariba}
\end{figure}

La Figura \ref{fig:captura-ses-ariba} presenta el módulo de SES/Ariba, orientado al control de hojas de entrada de servicio, referencias, soportes y aprobación antes de facturar. Esta pantalla se relaciona con los Pasos 10 y 11 del flujo. Su aporte operativo consiste en mantener trazabilidad de los soportes que habilitan la facturación, sin afirmar integración directa con SAP Ariba ni emisión automática de SES.

\begin{figure}[H]
\centering
\AppScreenshot[width=0.96\textwidth]{imagenes/captura_cierre_administrativo.png}
\caption{Resumen de cierre administrativo por actas, facturas y pagos. Fuente: captura propia del aplicativo CERMONT.}
\label{fig:captura-cierre-administrativo}
\end{figure}

La Figura \ref{fig:captura-cierre-administrativo} evidencia la vista consolidada de cierre administrativo. La pantalla agrupa actas, facturas y pagos, elementos vinculados con los Pasos 8 a 14. Su función es ofrecer una lectura de estado para el tramo final del proceso, donde una orden ejecutada debe convertirse en soporte aprobado, factura radicada y pago registrado.

\begin{figure}[H]
\centering
\AppScreenshot[width=0.96\textwidth]{imagenes/captura_motor_costos.png}
\caption{Motor de costos y consulta por orden. Fuente: captura propia del aplicativo CERMONT.}
\label{fig:captura-motor-costos}
\end{figure}

La Figura \ref{fig:captura-motor-costos} muestra la pantalla del motor de costos. Esta funcionalidad es transversal a los Pasos 3 a 14, porque conecta la propuesta económica, la planeación de recursos, el consumo real durante la ejecución y el seguimiento financiero del cierre. En esta versión, la captura evidencia el diseño del módulo y su interfaz de consulta; las métricas de desviación de costos quedan pendientes por anexar evidencia operativa con datos reales.

\begin{figure}[H]
\centering
\AppScreenshot[width=0.96\textwidth]{imagenes/captura_inventario_recursos.png}
\caption{Inventario y recursos asociados a la operación. Fuente: captura propia del aplicativo CERMONT.}
\label{fig:captura-inventario-recursos}
\end{figure}

La Figura \ref{fig:captura-inventario-recursos} corresponde al inventario de recursos. El módulo permite consultar herramientas, materiales, equipos y personal disponible para planeación y ejecución. Se relaciona principalmente con los Pasos 5 y 6, ya que la preparación correcta de recursos puede mitigar el riesgo de iniciar actividades sin los elementos necesarios. La captura evidencia la existencia de una vista de control, aunque los saldos y consumos reales deben documentarse con datos anexos de la empresa.

\begin{figure}[H]
\centering
\AppScreenshot[width=0.96\textwidth]{imagenes/captura_activos.png}
\caption{Vista de activos e intervenciones asociadas. Fuente: captura propia del aplicativo CERMONT.}
\label{fig:captura-activos}
\end{figure}

La Figura \ref{fig:captura-activos} muestra el módulo de activos. Esta pantalla permite relacionar intervenciones con elementos físicos o sistemas atendidos por CERMONT S.A.S. En el flujo operativo, los activos ayudan a contextualizar solicitudes, visitas, órdenes y evidencias, especialmente cuando una actividad técnica se repite sobre equipos, instalaciones o componentes que requieren historial.

\begin{figure}[H]
\centering
\AppScreenshot[width=0.96\textwidth]{imagenes/captura_perfil_usuario.png}
\caption{Perfil de usuario y cambio de credenciales. Fuente: captura propia del aplicativo CERMONT.}
\label{fig:captura-perfil-usuario}
\end{figure}

La Figura \ref{fig:captura-perfil-usuario} documenta la vista de perfil. Aunque no pertenece a un paso operativo específico, se vincula con la seguridad transversal del sistema, pues permite administrar información básica del usuario y cambio de contraseña. Esta evidencia complementa el modelo de autenticación y control de acceso descrito en la arquitectura, necesario para proteger acciones críticas dentro del flujo de 14 pasos.

\section{Arquitectura del software y estructura del monorepo}

Para soportar de forma robusta las exigencias del sistema, el proyecto se estructura bajo un esquema de monorepo gestionado mediante \textbf{npm workspaces} y optimizado con \textbf{Turborepo} para la orquestación de tareas de compilación, formateo y testing. La organización de directorios en el repositorio de producción es la siguiente:
\begin{itemize}
    \item \texttt{packages/shared-types}: Contiene la definición única de los contratos de datos de la aplicación. Utiliza la librería \textbf{Zod} para declarar esquemas de validación y exporta tipos TypeScript estrictos, actuando como la única fuente de verdad transaccional del proyecto.
    \item \texttt{backend}: API REST de alto rendimiento implementada en \textbf{Express 5.2.1} utilizando TypeScript estricto. Implementa la lógica de negocio pura, persistencia documental mediante Mongoose y la generación de documentos PDF con pdf-lib.
    \item \texttt{frontend}: Interfaz web moderna basada en el framework de producción \textbf{Next.js 16.2.4} con App Router, React 19 y Tailwind CSS para diseño visual cohesivo y mobile-first.
\end{itemize}

Esta separación física y conceptual de componentes y la centralización de tipos en el workspace compartido previene desviaciones y fallos de integración comunes en el desarrollo paralelo.

\begin{figure}[!htbp]
\centering
\resizebox{0.95\textwidth}{!}{%
\begin{tikzpicture}[node distance=0.9cm, every node/.style={font=\scriptsize}, layer/.style={draw, rounded corners, align=center, minimum width=4.2cm, minimum height=0.8cm, fill=blue!6}]
\node[layer] (ui) {Frontend\\Next.js / React / UI modular};
\node[layer, below=of ui] (contracts) {Contratos compartidos\\Zod / tipos / RBAC};
\node[layer, below=of contracts] (api) {Backend API\\Express / servicios de dominio};
\node[layer, below=of api] (data) {Persistencia\\MongoDB / Mongoose / archivos};
\draw[-{Latex[length=2mm]}] (ui)--(contracts); \draw[-{Latex[length=2mm]}] (contracts)--(api); \draw[-{Latex[length=2mm]}] (api)--(data);
\node[draw, rounded corners, fill=green!8, align=center, right=1.1cm of contracts, minimum width=3.4cm] {Auditoría\\validación\\permisos};
\draw[-{Latex[length=2mm]}, dashed] (contracts)--(api.east);
\end{tikzpicture}
%
}
\caption{Arquitectura por capas y flujo de integración del monorepo. Fuente: Elaboración propia.}
\label{fig:arquitectura-software}
\end{figure}

\section{Justificación del stack tecnológico seleccionado}
\label{sec:stack-justificado}

La selección del stack tecnológico se justifica a partir de necesidades verificables del proyecto y de la evidencia disponible en el repositorio. No se trata únicamente de un conjunto de herramientas contemporáneas, sino de una combinación orientada a resolver trazabilidad documental, heterogeneidad de formularios, control de acceso y operación con conectividad variable.

\begin{table}[!htbp]
\centering
\caption{Evidencia técnica del stack en el repositorio. Fuente: elaboración propia.}
\label{tab:justificacion-stack}
\footnotesize
\begin{adjustbox}{max width=\textwidth}
\begin{tabular}{|p{3.0cm}|p{4.5cm}|p{4.2cm}|p{3.0cm}|}
\hline\ClassicTableHeader
\textbf{Componente} & \textbf{Evidencia en el proyecto} & \textbf{Función dentro del flujo} & \textbf{Límite declarado}\\
\hline
Monorepo npm & Carpetas \texttt{backend}, \texttt{frontend} y \texttt{packages}. & Comparte contratos y roles entre módulos del flujo. & Requiere control de cambios transversales.\\
Express y rutas API & Controladores, servicios y middlewares backend. & Expone operaciones para órdenes, evidencias, actas, SES y pagos. & La disciplina de capas debe mantenerse manualmente.\\
MongoDB y Mongoose & Modelos documentales y subesquemas. & Conserva datos heterogéneos asociados a órdenes y soportes. & Requiere reglas explícitas de consistencia.\\
Next.js y React & Vistas del panel, formularios y rutas protegidas. & Presenta cada módulo del flujo a los usuarios autorizados. & Debe separar componentes de servidor y cliente.\\
Zod y shared-types & Esquemas compartidos y contratos exportados. & Valida payloads usados por frontend y backend. & No reemplaza reglas de negocio del servicio.\\
PWA y sincronización & Serwist/Turbopack, Dexie, TanStack Query persistente, cola offline, estado de conexión con ping real y resultados por ítem. & Permite continuidad parcial en actividades de campo y consulta offline de rutas previamente calentadas. & Cobertura offline completa por formulario: pendiente por anexar evidencia E2E para cada mutación crítica.\\
Generación documental & Servicios de PDF y módulo de documentos. & Apoya informes, actas y soportes de cierre. & Depende de plantillas y datos completos.\\
\hline
\end{tabular}
\end{adjustbox}
\end{table}

\section{El contrato compartido: validación bajo enfoque de confianza cero (Zero Trust) con Zod}

Un pilar de seguridad y robustez en la plataforma es la reducción de validaciones redundantes o divergentes entre capas. Mediante el uso de \texttt{@cermont/shared-types}, un único esquema Zod (ej. \texttt{WorkOrderSchema}, \texttt{ProposalSchema}, \texttt{EvidenceSchema}) es importado tanto por el frontend como por el backend:
\begin{enumerate}
    \item \textbf{Frontend (Validación de UX)}: Los esquemas de Zod se acoplan directamente a las vistas del lado del cliente mediante la integración de \textbf{React Hook Form}. Esto proporciona retroalimentación inmediata sobre el formato, campos obligatorios y tipos de datos sin realizar peticiones de red, lo cual favorece una captura más clara para el técnico en campo.
    \item \textbf{Backend (Validación de Seguridad)}: Antes de realizar cualquier operación de persistencia o procesamiento en la base de datos documental, un middleware centralizado en Express ejecuta un análisis sintáctico (\texttt{safeParse}) del payload utilizando el mismo esquema de Zod. Si el cliente ha enviado datos alterados o incompletos, la petición se rechaza de inmediato con un código HTTP 400 Bad Request, implementando un perímetro estricto de confianza cero (Zero Trust).
\end{enumerate}

\section{Estructura técnica del Backend en Express 5.2.1}

La arquitectura del backend se organiza bajo un patrón clásico MVC (Model-View-Controller) modificado para servicios asíncronos y desacoplados. La estructura interna de cada caso de uso sigue un flujo lineal de responsabilidades aisladas:
\begin{equation}
\text{Router} \longrightarrow \text{Controller} \longrightarrow \text{Service} \longrightarrow \text{Model (Mongoose)}
\end{equation}

\begin{itemize}
    \item \textbf{Router}: Define la ruta del endpoint HTTP, los middlewares de seguridad (RBAC) y la validación de payload Zod.
    \item \textbf{Controller}: Extrae los parámetros de la petición HTTP, delega la ejecución de la regla de negocio al Service correspondiente y formatea la respuesta en un envoltorio estandarizado:
    \begin{lstlisting}[caption={Estructura estándar de respuesta API}, label={lst:api-response}]
{ success: boolean, data: T | null, error: string | null, message: string }
\end{lstlisting}
    Cabe destacar que gracias a la migración a \textbf{Express 5.2.1}, no se requiere envolver los métodos de los controladores en bloques \texttt{try/catch} repetitivos o funciones \texttt{asyncHandler}. Express 5 propaga de forma nativa los errores asíncronos lanzados en las promesas directamente al middleware global de gestión de errores del servidor.
    \item \textbf{Service}: Clase o módulo puro de TypeScript que encapsula la lógica operativa pura de negocio. Tiene prohibido importar o interactuar con los objetos de red de Express (\texttt{Request} o \texttt{Response}), facilitando su testabilidad unitaria aislada.
    \item \textbf{Model (Mongoose 9.x)}: Capa de persistencia que interactúa de manera directa con MongoDB, aplicando modelos Mongoose, validaciones de esquema y referencias mediante \texttt{ObjectId} cuando una entidad requiere identidad independiente.
\end{itemize}

\section{Diseño del Frontend y el Perímetro de Seguridad en Next.js 16}

La aplicación frontend implementa una estructura de interfaz de usuario mobile-first que escala progresivamente para dispositivos de escritorio utilizando la directiva de Tailwind CSS de forma nativa. Para garantizar la consistencia, la arquitectura web adopta los siguientes patrones de desarrollo:
\begin{itemize}
    \item \textbf{TanStack Query v5 (Manejo del Estado del Servidor)}: Todas las operaciones de obtención y mutación de datos de la API se realizan a través de hooks personalizados que delegan el almacenamiento en caché, la invalidación de datos, los reintentos automáticos y el estado de carga (skeletons) a TanStack Query, aislando la red de los componentes React.
    \item \textbf{Zustand v5 (Manejo del Estado Local)}: Se utiliza para almacenar el estado volátil de la interfaz (ej. barra lateral colapsada) y la persistencia segura del token de sesión a través del almacén global \texttt{auth.store.ts}.
    \item \textbf{proxy.ts como perímetro de seguridad}: En cumplimiento con la arquitectura definida para el proyecto, el aplicativo implementa un archivo centralizado \texttt{proxy.ts} que intercepta peticiones entrantes del lado del cliente, verifica la existencia del token JWT almacenado en cookies HttpOnly y evalúa la matriz de permisos RBAC antes de renderizar la página o delegar la llamada a la API del backend. Este control fortalece el aislamiento de los módulos del panel de administración.
\end{itemize}

\begin{figure}[!htbp]
\centering
\resizebox{0.95\textwidth}{!}{%
\begin{tikzpicture}[
  node distance=0.78cm,
  every node/.style={font=\tiny},
  ent/.style={draw=institucionalBlue, rounded corners=3pt, fill=institucionalGrey!35, align=center, minimum width=2.12cm, minimum height=0.62cm, line width=0.65pt},
  entred/.style={draw=institucionalRed, rounded corners=3pt, fill=institucionalRed!8, align=center, minimum width=2.12cm, minimum height=0.62cm, line width=0.75pt},
  link/.style={-{Latex[length=1.6mm]}, draw=institucionalBlue, line width=0.65pt},
  linkred/.style={-{Latex[length=1.6mm]}, draw=institucionalRed, line width=0.75pt}
]
\node[ent] (wr) {WorkRequest};
\node[ent, right=of wr] (pr) {Proposal};
\node[ent, right=of pr] (po) {PurchaseOrder};
\node[entred, below=of pr] (wo) {WorkOrder};
\node[ent, below left=of wo] (plan) {PlanningPacket};
\node[ent, below=of wo] (exe) {ExecutionSession};
\node[ent, below right=of wo] (ev) {Evidence};
\node[ent, below=of exe] (rep) {TechnicalReport};
\node[ent, right=of rep] (acta) {DeliveryRecord};
\node[ent, right=of acta] (ses) {SES};
\node[ent, right=of ses] (inv) {Invoice};
\node[entred, right=of inv] (pay) {Payment};

\draw[link] (wr)--(pr);
\draw[link] (pr)--(po);
\draw[linkred] (pr)--(wo);
\draw[linkred] (po)--(wo);
\draw[linkred] (wo)--(plan);
\draw[linkred] (wo)--(exe);
\draw[link] (exe)--(ev);
\draw[link] (exe)--(rep);
\draw[link] (rep)--(acta);
\draw[linkred] (acta)--(ses);
\draw[linkred] (ses)--(inv);
\draw[linkred] (inv)--(pay);
\end{tikzpicture}
%
}
\caption{Diagrama de base documental y relaciones lógicas. Fuente: Elaboración propia.}
\label{fig:modelo-datos}
\end{figure}

\section{Operación offline, persistencia local y sincronización diferida}

Para atender escenarios de conectividad móvil limitada, el aplicativo implementa una estrategia de persistencia local basada en Serwist/Turbopack, Service Worker, Dexie/IndexedDB, TanStack Query persistente y una cola de sincronización diferida. La evidencia del repositorio muestra que esta capacidad existe, pero su alcance debe declararse como parcial: permite navegación offline posterior a una carga inicial, fallback \path{/~offline} para rutas no visitadas, almacenamiento local de operaciones pendientes, snapshots locales de listados operativos y sincronización con resultados por ítem; sin embargo, no convierte en definitivos los procesos externos de SES, facturación o pago sin confirmación del backend.

\begin{enumerate}
    \item \textbf{Service Worker de producción}: El archivo fuente \path{frontend/src/app/sw.ts} se compila con \texttt{@serwist/turbopack} y se sirve desde \path{/serwist/sw.js}. El worker precachea \path{/~offline}, aplica \texttt{CacheFirst} a \path{/_next/static/*}, \texttt{StaleWhileRevalidate} a imágenes y \texttt{NetworkFirst} a documentos internos visitados.
    \item \textbf{Detección de desconexión}: La interfaz no depende exclusivamente de \texttt{navigator.onLine}; utiliza señales del navegador como indicio inicial y confirma conectividad con peticiones \texttt{HEAD} a \path{/api/backend/health} y, como respaldo, a \path{/serwist/sw.js}, con tiempo de espera y retroceso progresivo.
    \item \textbf{Persistencia local}: Los hooks de operación offline almacenan en la cola local la mutación pendiente junto con su \texttt{idempotencyKey} o identificador de mutación local. Además, los listados de casos de servicio, solicitudes, visitas técnicas y plantillas documentales conservan snapshots en IndexedDB para consulta posterior.
    \item \textbf{Actualización visible del estado}: La interfaz informa al usuario que la acción quedó almacenada localmente o que se están mostrando datos guardados localmente, según el módulo y el estado de sincronización.
    \item \textbf{Sincronización diferida}: Al restablecerse la red, el administrador de sincronización reenvía las operaciones compatibles al backend y marca cada ítem como sincronizado, fallido o en conflicto según la respuesta recibida.
    \item \textbf{Límite actual}: La captura offline completa debe validarse por formulario mediante pruebas E2E. Por tanto, el libro clasifica esta capacidad como parcialmente implementada y no como operación offline total del flujo de 14 pasos.
\end{enumerate}

\begin{figure}[!htbp]
\centering
\resizebox{0.95\textwidth}{!}{%
\begin{tikzpicture}[
  node distance=0.9cm,
  every node/.style={font=\scriptsize},
  box/.style={draw=institucionalBlue, rounded corners=3pt, align=center, minimum width=2.45cm, minimum height=0.78cm, fill=institucionalGrey!35, line width=0.75pt},
  limit/.style={draw=institucionalRed, rounded corners=3pt, align=center, text width=4.1cm, minimum height=0.58cm, fill=institucionalRed!8, font=\tiny, line width=0.75pt},
  arrow/.style={-{Latex[length=2mm]}, draw=institucionalBlue, line width=0.8pt},
  warn/.style={-{Latex[length=2mm]}, draw=institucionalRed, line width=0.75pt}
]
\node[box] (cap) {Captura en campo};
\node[box, right=of cap] (local) {IndexedDB\\cola local};
\node[box, right=of local] (net) {Detección\\de red};
\node[box, right=of net] (api) {API backend};
\node[box, right=of api] (ok) {Confirmación\\y auditoría};
\foreach \a/\b in {cap/local,local/net,net/api,api/ok}{\draw[arrow] (\a)--(\b);}

\node[limit, below=0.62cm of local] (note) {La sincronización completa de archivos binarios se declara como alcance parcial};
\draw[warn] (local) -- (note);
\end{tikzpicture}
%
}
\caption[Secuencia de sincronización offline]{Secuencia de sincronización offline: captura local, almacenamiento en IndexedDB, envío diferido y confirmación. Fuente: Elaboración propia.}
\label{fig:secuencia-offline-sync}
\end{figure}

\begin{figure}[!htbp]
\centering
\resizebox{0.95\textwidth}{!}{%
\begin{tikzpicture}[
  node distance=0.75cm,
  every node/.style={font=\scriptsize},
  box/.style={draw=institucionalBlue, rounded corners=3pt, align=center, minimum width=2.18cm, minimum height=0.72cm, fill=institucionalGrey!35, line width=0.75pt},
  review/.style={draw=institucionalRed, rounded corners=3pt, align=center, minimum width=2.18cm, minimum height=0.72cm, fill=institucionalRed!8, line width=0.85pt},
  arrow/.style={-{Latex[length=2mm]}, draw=institucionalBlue, line width=0.8pt},
  arrowred/.style={-{Latex[length=2mm]}, draw=institucionalRed, line width=0.8pt}
]
\node[box] (up) {Documento\\PDF / Word / foto};
\node[box, right=of up] (imp) {Importación};
\node[box, right=of imp] (map) {Mapeo de campos};
\node[review, right=of map] (rev) {Revisión humana};
\node[box, right=of rev] (form) {Formulario digital};
\node[box, right=of form] (pdf) {Documento generado};

\foreach \a/\b in {up/imp,imp/map,map/rev,rev/form,form/pdf}{\draw[arrow] (\a)--(\b);}
\draw[arrowred] (rev.south) -- ++(0,-0.42) node[below, font=\tiny, align=center] {control académico y operativo};
\end{tikzpicture}
%
}
\caption{Flujo propuesto para extracción estructurada de documentos. Fuente: Elaboración propia.}
\label{fig:pipeline-formularios}
\end{figure}

% =====================================================================
\section{Arquitectura de seguridad en profundidad}
\label{sec:arquitectura-seguridad}

La protección de los datos operativos y financieros que transitan por el aplicativo se diseñó bajo el principio de defensa en profundidad (\textit{Defense in Depth}), que establece que ninguna capa individual de seguridad debe considerarse suficiente por sí sola. El sistema implementa seis capas de protección dispuestas desde el perímetro externo hasta el núcleo de persistencia.

\subsection{Primera capa: Perímetro de red (Proxy y CORS)}

El archivo \texttt{proxy.ts} en el frontend actúa como la primera barrera de seguridad. Intercepta todas las peticiones entrantes antes de que alcancen el enrutador de Next.js y aplica tres verificaciones: (a) redirección de rutas públicas a la página de inicio de sesión cuando el usuario no está autenticado; (b) validación de la existencia y vigencia del token JWT en la cookie HttpOnly; y (c) verificación de que la ruta solicitada pertenece al conjunto de rutas permitidas para el rol del usuario, utilizando las definiciones centralizadas del paquete \texttt{@cermont/domain}. En el backend, el middleware CORS restringe las peticiones exclusivamente al origen del frontend, configurado mediante variables de entorno.

\subsection{Segunda capa: Autenticación sin estado (JWT)}

Tras la validación inicial de credenciales contra el hash Bcrypt almacenado, el backend emite dos artefactos: un token de acceso JWT con tiempo de vida corto (24 horas) y un token de actualización (\textit{refresh token}) almacenado en una cookie HttpOnly con los flags \texttt{Secure}, \texttt{SameSite=Strict} y \texttt{Path=/api/auth}. El token de acceso viaja exclusivamente en memoria del lado del cliente (almacén Zustand), lo que lo hace inmune a ataques XSS que intenten leer \texttt{localStorage}. El middleware \texttt{authenticate} en el backend verifica la firma criptográfica, la expiración y la integridad del token en cada petición a endpoints protegidos.

\subsection{Tercera capa: Autorización basada en roles (RBAC)}

Una vez autenticado el usuario, el middleware \texttt{authorize} consulta el rol asignado y lo contrasta contra la matriz de permisos del módulo solicitado. Por ejemplo, un usuario con rol \texttt{tecnico} puede consultar y actualizar órdenes de trabajo que le han sido asignadas, pero no puede aprobar propuestas ni crear usuarios. Esta verificación se realiza en el backend para cada endpoint protegido; el frontend únicamente oculta o deshabilita elementos de interfaz como mejora de experiencia de usuario, nunca como mecanismo de seguridad.

\subsection{Cuarta capa: Validación de entradas (Zero Trust)}

Todo payload recibido por el backend —sin excepción— es validado mediante un esquema Zod antes de alcanzar la lógica de negocio. Esta validación verifica tipos de datos, formatos, longitudes, rangos numéricos, valores permitidos y campos obligatorios. Un payload que no supera la validación recibe un error HTTP 400 con un código de error tipado que describe exactamente qué campo falló y por qué. Esta capa protege contra inyección NoSQL, desbordamiento de buffer, y datos malformados intencional o accidentalmente.

\subsection{Quinta capa: Registro de auditoría inmutable}

Cada acción crítica —creación de orden de trabajo, cambio de estado, aprobación de propuesta, generación de SES, emisión de factura, cambio de rol de usuario— genera un evento de auditoría inmutable que registra: tipo de evento, identificador del usuario que ejecutó la acción, entidad afectada, estado anterior, estado nuevo, marca de tiempo y metadatos relevantes. Estos registros no pueden ser modificados ni eliminados por ningún rol del sistema, y proporcionan la trazabilidad requerida para auditorías internas y externas.

\subsection{Sexta capa: Protección de datos en reposo y en tránsito}

Las contraseñas se almacenan exclusivamente como hashes Bcrypt con factor de costo configurable. En un despliegue productivo, las comunicaciones entre el frontend y el backend deben protegerse mediante HTTPS/TLS configurado en el servidor o proxy correspondiente. Los archivos cargados (evidencias fotográficas, documentos PDF) se procesan con \texttt{sharp} para eliminar metadatos EXIF potencialmente sensibles (como coordenadas GPS de las fotografías) antes de su almacenamiento, conservando únicamente la información de geolocalización cuando es requerida por el flujo operativo y se registra de forma explícita.

\section{Arquitectura de módulos del sistema: los catorce pasos del flujo CERMONT}
\label{sec:modulos-flujo}

El aplicativo implementa el flujo operativo completo de CERMONT S.A.S., compuesto por catorce pasos que abarcan desde la solicitud inicial del cliente hasta el pago final. A continuación se describen los primeros siete pasos, correspondientes al ciclo operativo; los siete pasos restantes abarcan el cierre administrativo y se detallan en la subsección siguiente.

\subsection{Paso 1: Solicitud de trabajo (Work Request)}

\textbf{Entidad:} WorkRequest. \textbf{Estados:} open $\rightarrow$ assigned $\rightarrow$ in\_progress $\rightarrow$ completed $\rightarrow$ cancelled. \textbf{API:} \path{GET/POST/PATCH /api/work-requests}. \textbf{RBAC:} gerente, residente, HES (CRUD completo); cliente (crear y leer propias).

La solicitud de trabajo constituye el punto de entrada al sistema. Un cliente registra una necesidad de servicio especificando tipo de trabajo, ubicación, urgencia y documentos adjuntos. El sistema notifica al residente asignado, quien evalúa si la solicitud requiere visita técnica o puede proceder directamente a propuesta.

\subsection{Paso 2: Visita técnica (Site Visit)}

\textbf{Entidad:} SiteVisit. \textbf{Estados:} pending $\rightarrow$ scheduled $\rightarrow$ in\_progress $\rightarrow$ completed $\rightarrow$ cancelled. \textbf{API:} \path{GET/POST /api/work-requests/:id/visits}. \textbf{Offline:} consulta de listado con snapshot local; la captura completa sin conexión debe validarse por formulario.

Cuando la complejidad del servicio lo requiere, se agenda una visita al sitio. El técnico asignado registra mediciones, condiciones del lugar, requisitos de seguridad y evidencia fotográfica preliminar. Estos datos pueden alimentar la propuesta económica y reducir recaptura manual cuando el flujo se usa de extremo a extremo.

\subsection{Paso 3: Propuesta económica (Proposal)}

\textbf{Entidad:} Proposal. \textbf{Estados:} draft $\rightarrow$ sent $\rightarrow$ approved $\rightarrow$ rejected $\rightarrow$ expired. \textbf{API:} \path{GET/POST /api/proposals}.

A partir de los datos de la solicitud y, cuando existe, de la visita técnica, el sistema permite generar una propuesta económica estructurada que incluye: alcance del servicio, recursos estimados (personal, materiales, equipos), cronograma, condiciones comerciales y valor total. La propuesta se exporta como PDF para envío al cliente.

\subsection{Paso 4: Aprobación con orden de compra (Purchase Order)}

\textbf{Entidad:} PurchaseOrder. \textbf{Estados:} pending $\rightarrow$ received $\rightarrow$ approved $\rightarrow$ rejected. \textbf{API:} \path{POST /api/proposals/:id/po}.

El cliente aprueba la propuesta mediante una orden de compra (PO) que se adjunta al sistema. La recepción de la PO habilita la creación de la orden de trabajo interna y autoriza el inicio de la planeación de recursos. Este paso actúa como compuerta contractual: sin PO aprobada, no se puede crear una orden de trabajo.

\subsection{Entidad transversal: Orden de trabajo (WorkOrder / ServiceCase)}

\textbf{Entidad:} WorkOrder. \textbf{Rol en el sistema:} entidad raíz de trazabilidad. \textbf{API:} \path{GET/POST/PATCH /api/orders}.

La orden de trabajo no se presenta como un paso independiente del flujo original, sino como la entidad digital que articula todos los pasos del proceso. Agrupa la solicitud, la visita técnica, la propuesta, la orden de compra, la planeación, la ejecución, las evidencias, el informe técnico, el acta, la SES, el seguimiento de factura y el registro de pago. Su máquina de estados refleja el avance global del servicio y determina qué módulos se habilitan en cada momento.

\begin{figure}[!htbp]
\centering
\resizebox{0.95\textwidth}{!}{%
\begin{tikzpicture}[
  every node/.style={font=\tiny},
  state/.style={draw=institucionalBlue, rounded corners=3pt, fill=institucionalGrey!35, align=center, minimum width=1.86cm, minimum height=0.56cm, line width=0.65pt},
  critical/.style={draw=institucionalRed, rounded corners=3pt, fill=institucionalRed!8, align=center, minimum width=1.86cm, minimum height=0.56cm, line width=0.75pt},
  arrow/.style={-{Latex[length=1.35mm]}, draw=institucionalBlue, line width=0.62pt},
  arrowred/.style={-{Latex[length=1.35mm]}, draw=institucionalRed, line width=0.75pt}
]
\foreach \i/\txt in {0/Creada,1/Planeación,2/Lista,3/Ejecución,4/Terminada,5/Informe,6/Acta,7/SES,8/Factura,9/Pago,10/Cerrada}{
  \ifnum\i>6
    \node[critical] (s\i) at (\i*1.76,0) {\txt};
  \else
    \node[state] (s\i) at (\i*1.76,0) {\txt};
  \fi
}
\foreach \i in {0,...,6}{\pgfmathtruncatemacro{\j}{\i+1}\draw[arrow] (s\i)--(s\j);}
\foreach \i in {7,...,9}{\pgfmathtruncatemacro{\j}{\i+1}\draw[arrowred] (s\i)--(s\j);}
\node[below=0.72cm of s5, align=center, font=\scriptsize, text=institucionalBlue] {Cada transición requiere precondiciones, permisos y evidencia documental};
\end{tikzpicture}
%
}
\caption{Máquina de estados finita (FSM) que controla la trazabilidad operativa y los bloqueos de seguridad de la orden. Fuente: elaboración propia.}
\label{fig:fsm-15-estados}
\end{figure}

\subsection{Paso 5: Planeación de recursos (Planning Packet)}

\textbf{Entidad:} PlanningPacket. \textbf{Estados:} draft $\rightarrow$ ready $\rightarrow$ approved. \textbf{Relación con el problema:} corrige la planeación incompleta de herramientas, equipos, personal, EPP, permisos y documentos de apoyo.

En esta fase se asignan los recursos necesarios para la ejecución: personal técnico con certificaciones vigentes, herramientas y equipos del catálogo, materiales y consumibles, elementos de protección personal (EPP), vehículos y documentos requeridos (permisos de trabajo, AST, procedimientos, instructivos y formatos de tareas críticas). El sistema permite reutilizar kits típicos por tipo de actividad, disminuyendo la dependencia de memoria individual y evitando que cada orden se planifique completamente desde cero.

\subsection{Paso 6: Ejecución en campo (Execution Session)}

\textbf{Entidad:} ExecutionSession. \textbf{Estados:} scheduled $\rightarrow$ in\_progress $\rightarrow$ paused $\rightarrow$ completed $\rightarrow$ cancelled. \textbf{Offline:} Sí, prioritario.

Durante esta fase, los técnicos en campo confirman el consumo de materiales y horas de trabajo, diligencian checklists digitales de seguridad e inspección, capturan evidencias fotográficas vinculadas a la orden y registran novedades, hallazgos o acciones correctivas. La captura offline evita que la ausencia de conectividad impida el registro de información crítica; la sincronización posterior se realiza cuando el dispositivo recupera conexión.

\subsection{Paso 7: Consolidación del informe técnico (Technical Report)}

\textbf{Entidad:} TechnicalReport. \textbf{Estado FSM:} \texttt{report\_pending} $\rightarrow$ \texttt{completed}. \textbf{API:} \path{POST /api/reports}.

Una vez finalizada la ejecución, el sistema consolida los datos capturados en campo, las listas de verificación, las observaciones y las evidencias fotográficas para generar un informe técnico preliminar. Este documento no sustituye la revisión profesional del responsable, sino que proporciona una base estructurada para la validación del residente o supervisor y puede reducir recaptura manual cuando los datos de origen están completos.

\subsection{Pasos 8 al 14: Cierre administrativo y seguimiento de recaudo}

Los pasos finales del flujo CERMONT buscan que la ejecución técnica se traduzca de forma efectiva en facturación y recaudo. El sistema implementa un módulo de cierre consolidado (\texttt{ClosureReport}) que actúa como un modelo de lectura (\textit{Read Model}) agregando información de cinco entidades independientes para proporcionar trazabilidad consolidada.

\subsubsection{Paso 8: Acta de Entrega Técnica (Delivery Record)}
\textbf{Entidad:} DeliveryRecord. \textbf{Estado FSM:} \texttt{completed} $\rightarrow$ \texttt{ready\_for\_invoicing}. \textbf{API:} \path{POST /api/delivery-records}.
Al finalizar el informe técnico, se genera el acta de entrega que detalla observaciones finales y recursos consumidos. El sistema valida que el informe técnico esté aprobado antes de permitir la creación del acta.

\subsubsection{Paso 9: Firma del Acta (Client Acceptance)}
\textbf{Estado FSM:} \texttt{ready\_for\_invoicing} $\rightarrow$ \texttt{acta\_signed}. \textbf{API:} \path{PATCH /api/delivery-records/:id/sign}.
El sistema registra la aceptación del cliente capturando el nombre del responsable, la fecha y el método de firma (manual con soporte cargado o digital de texto). Este paso habilita formalmente el proceso de cobro.

\subsubsection{Paso 10: Elaboración y envío de SES}
\textbf{Entidad:} ServiceEntrySheet (SES). \textbf{Estado FSM:} \texttt{acta\_signed} $\rightarrow$ \texttt{ses\_sent}. \textbf{API:} \path{POST /api/service-entry-sheets}.

El sistema registra el número de documento de Ariba, la fecha de radicación, el responsable interno y el soporte documental asociado. La integración directa con SAP Ariba se declara fuera del alcance de esta versión; por tanto, el aplicativo funciona como espejo interno de seguimiento para evitar pérdida de trazabilidad.

\subsubsection{Paso 11: Aprobación de SES}
\textbf{Entidad:} SESApproval. \textbf{Estado operativo:} seguimiento interno de aprobación. \textbf{API:} \path{PATCH /api/service-entry-sheets/:id}.

Cuando el cliente aprueba la SES, el responsable administrativo registra la fecha de aprobación y el soporte correspondiente. Si la SES es rechazada o requiere corrección, el sistema conserva el historial de observaciones para evitar que el proceso quede únicamente en correos o mensajes informales.

\subsubsection{Paso 12: Elaboración y envío de factura}
\textbf{Entidad:} InvoiceTracking. \textbf{Estado FSM:} \texttt{ses\_sent} $\rightarrow$ \texttt{invoice\_approved} cuando el cliente aprueba la factura. \textbf{API:} \path{POST /api/invoices}.

El aplicativo no emite factura electrónica DIAN ni reemplaza el software contable certificado de la empresa. Su función es registrar el seguimiento documental de la factura emitida externamente: número, fecha, valor, soporte PDF, responsable y estado de radicación.

\subsubsection{Paso 13: Aprobación de factura}
\textbf{Entidad:} InvoiceApproval. \textbf{Estado operativo:} aprobado, rechazado o en corrección. \textbf{API:} \path{PATCH /api/invoices/:id/status}.

El sistema permite registrar la aprobación o devolución de la factura por parte del cliente. Esta información alimenta el tablero administrativo y permite identificar órdenes que ya fueron ejecutadas técnicamente, pero permanecen bloqueadas por requisitos documentales o financieros.

\subsubsection{Paso 14: Registro de Pago y Cierre Definitivo}
\textbf{Entidad:} Payment. \textbf{Estado FSM:} \texttt{invoice\_approved} $\rightarrow$ \texttt{paid} $\rightarrow$ \texttt{closed}. \textbf{API:} \path{POST /api/payments}.
El flujo concluye con el registro interno del pago mediante referencia bancaria, fecha de abono y soporte documental. Una vez verificado el pago por el área administrativa, la orden de trabajo transiciona al estado terminal \texttt{closed}, marcando la finalización del ciclo de trazabilidad de 14 pasos.

\section{Arquitectura de la capa de persistencia}
\label{sec:persistencia}

La capa de persistencia del sistema se apoya en MongoDB 7.0, una base de datos orientada a documentos que almacena los datos en formato BSON (Binary JSON). Esta elección técnica responde a la naturaleza heterogénea de los formularios operativos de CERMONT S.A.S., donde cada tipo de servicio (inspección de líneas de vida, mantenimiento CCTV, planeación de obra) requiere capturar conjuntos de campos diferentes.

\subsection{Patrones de modelado}

El diseño del esquema de base de datos aplica dos patrones complementarios:

\textbf{Documentos embebidos (\textit{Embedding}).} Cuando los datos tienen un ciclo de vida acoplado y se acceden siempre en conjunto, se almacenan como subdocumentos dentro del documento principal. Por ejemplo, las líneas de recursos de un kit (\texttt{KitResourceLine}) se embeben dentro del documento del kit, ya que carecen de identidad independiente fuera de él. De igual forma, las respuestas a un checklist se embeben dentro de la sesión de ejecución.

\textbf{Referencias (\textit{Referencing}).} Cuando una entidad tiene identidad independiente y es referenciada desde múltiples documentos, se almacena en su propia colección y se referencia mediante \texttt{ObjectId}. Por ejemplo, un \texttt{Asset} (activo físico) es referenciado desde órdenes de trabajo, planes de mantenimiento y sesiones de ejecución.

\subsection{Estrategia de índices}

Cada colección define índices alineados con los patrones de consulta más frecuentes: búsqueda por código único, filtrado por estado, consulta por entidad relacionada (\texttt{serviceCaseId}, \texttt{workOrderId}), ordenamiento por fecha de creación y búsqueda por texto en campos descriptivos. Los índices compuestos cubren las consultas más comunes, como listar todas las órdenes de trabajo de un cliente filtrando por estado y ordenadas por fecha.

\subsection{Campos de auditoría}

Todos los documentos en la base de datos incluyen un conjunto estándar de campos de auditoría: \texttt{createdBy} (identificador del usuario que creó el registro), \texttt{createdAt} (marca de tiempo de creación), \texttt{updatedBy} y \texttt{updatedAt} (última modificación). Para las entidades que soportan archivado lógico, se añaden \texttt{archivedAt} y \texttt{archivedBy}. Estos campos, combinados con la colección de eventos de auditoría, proporcionan trazabilidad completa de cada cambio en el sistema.

% =====================================================================
\section{Integración del ecosistema de desarrollo}
\label{sec:ecosistema-desarrollo}

El desarrollo del aplicativo se apoyó en un ecosistema de herramientas que automatizaron la verificación de calidad del código en cada iteración. A continuación se describe el rol de cada herramienta en el flujo de trabajo.

\textbf{TypeScript en modo estricto.} Todo el código del monorepo —backend, frontend y paquetes compartidos— se escribe en TypeScript con la configuración \texttt{strict: true}. Esto habilita verificaciones como \texttt{strictNullChecks}, \texttt{noImplicitAny} y \texttt{strictFunctionTypes}, que previenen categorías enteras de errores en tiempo de compilación. El comando \texttt{tsc --noEmit} ejecutado desde la raíz del monorepo verifica la corrección de tipos en los tres workspaces simultáneamente, sin necesidad de compilar a JavaScript.

\textbf{Biome para linting y formateo.} En lugar de la combinación tradicional ESLint + Prettier, el proyecto utiliza Biome, una herramienta unificada escrita en Rust que proporciona análisis estático y formateo de código. Biome verifica el cumplimiento de las reglas de estilo definidas para el proyecto, incluyendo la prohibición de \texttt{console.log}, \texttt{debugger}, variables no utilizadas e importaciones no usadas.

\textbf{Vitest para pruebas unitarias e integración.} Vitest, compatible con el ecosistema Vite, ejecuta las pruebas del proyecto con soporte nativo para TypeScript, sin requerir configuración adicional de transpilación. Su integración con Supertest permite probar los endpoints HTTP del backend realizando peticiones reales a una instancia de Express, mientras que la base de datos se aísla mediante una instancia de MongoDB en memoria (conexión a \texttt{mongodb://127.0.0.1:27017/cermont\_test}) que se crea antes de cada suite de pruebas y se destruye al finalizar.

\textbf{React Doctor para auditoría de componentes.} La herramienta React Doctor analiza el código del frontend en busca de problemas comunes: componentes excesivamente grandes, exports no utilizados, uso de \texttt{useSearchParams} o \texttt{useRouter} sin envoltura Suspense, creación de objetos \texttt{new Date()} o llamadas a \texttt{Math.random()} durante el renderizado (causa potencial de errores de hidratación), y problemas de accesibilidad. Su ejecución periódica ayuda a detectar problemas tempranos; la evidencia de cada corrida debe anexarse al informe de validación.

\textbf{Husky y lint-staged para ganchos de pre-commit.} Para prevenir que código con errores llegue al repositorio, se configuraron ganchos de Git mediante Husky que ejecutan automáticamente \texttt{lint-staged} sobre los archivos modificados antes de cada commit. Esto asegura que solo el código que pasa las verificaciones de Biome y TypeScript sea incorporado al historial del proyecto.

\section{Experiencia de desarrollo y curva de aprendizaje}
\label{sec:experiencia-desarrollo}

El desarrollo del aplicativo, realizado por un único ingeniero como parte de la práctica empresarial, permite extraer observaciones sobre la viabilidad de construir sistemas empresariales modernos con recursos limitados. El stack tecnológico seleccionado —Next.js, Express, MongoDB, TypeScript y Tailwind CSS— resultó adecuado para un desarrollador con formación en Ingeniería Electrónica. La consistencia del lenguaje en frontend y backend simplificó el cambio entre capas, mientras que la organización del monorepo con tipos compartidos ayudó a controlar definiciones duplicadas.

La principal dificultad técnica encontrada durante el desarrollo fue la implementación de la sincronización offline con detección de conflictos. El diseño de una cola de operaciones en IndexedDB que mantuviera la idempotencia de las mutaciones y manejara correctamente los casos de borde (aplicación cerrada durante la sincronización, múltiples dispositivos operando sobre la misma orden, pérdida de conectividad a mitad de una sincronización) requirió iteraciones adicionales de diseño y prueba. Esta experiencia es consistente con la literatura sobre sistemas distribuidos: la gestión de estados offline es uno de los problemas más complejos en el desarrollo de aplicaciones web modernas \cite{richards2020}.

La adopción temprana de las compuertas de calidad (typecheck, lint, test, build) desde la primera iteración resultó ser una decisión acertada. Aunque inicialmente añadió fricción al flujo de desarrollo, la inversión se recuperó en iteraciones posteriores, cuando las compuertas detectaron regresiones que habrían sido difíciles de identificar manualmente en un código base creciente.

\section{Comparativa tecnológica: stack implementado vs. alternativas evaluadas}
\label{sec:comparativa-stack}

La Tabla \ref{tab:comparativa-stack} presenta una comparación resumida entre el stack tecnológico implementado y las principales alternativas evaluadas durante la fase de diseño, indicando para cada caso el criterio que motivó la decisión final.

\begin{table}[!htbp]
\centering
\caption{Comparativa del stack implementado frente a alternativas evaluadas. Fuente: Elaboración propia.}
\label{tab:comparativa-stack}
\small
\def\StackRow#1#2#3#4{\parbox[t]{2.5cm}{\raggedright #1} & \parbox[t]{3.0cm}{\raggedright #2} & \parbox[t]{3.0cm}{\raggedright #3} & \parbox[t]{3.3cm}{\raggedright #4}\\\hline}
\begin{tabular}{|l|l|l|l|}
\hline\ClassicTableHeader
\textbf{Capa} & \textbf{Stack implementado} & \textbf{Alternativa evaluada} & \textbf{Criterio de decisión} \\
\hline
\StackRow{Frontend}{Next.js 16 + React 19 + Tailwind CSS 4}{Vite + React SPA}{Server Components para vistas de datos; ecosistema unificado de enrutamiento, API routes y PWA.}
\StackRow{Backend}{Express 5.2.1 + TypeScript}{NestJS}{Simplicidad, menor curva de aprendizaje, propagación nativa de errores asíncronos en Express 5.}
\StackRow{Base de datos}{MongoDB 7.0 + Mongoose 9.x}{PostgreSQL + Prisma}{Flexibilidad de esquema para formularios heterogéneos; sin migraciones para cada nuevo tipo de servicio.}
\StackRow{Validación}{Zod 4.x (compartido)}{Joi}{Inferencia de tipos TypeScript desde esquemas; ecosistema unificado frontend/backend.}
\StackRow{Autenticación}{JWT + cookies HttpOnly}{Auth.js / NextAuth}{Operación sin estado compatible con modo offline; sin dependencia de base de datos de sesiones.}
\StackRow{Pruebas}{Vitest + Supertest}{Jest}{Velocidad de ejecución; integración nativa con Vite y TypeScript; configuración mínima.}
\hline
\end{tabular}
\end{table}

Esta comparativa no pretende establecer una superioridad absoluta del stack implementado, sino documentar que cada decisión tecnológica fue el resultado de un análisis contextual que consideró las restricciones específicas del proyecto CERMONT: operación offline en zonas remotas, diversidad de formatos operativos, equipo de desarrollo unipersonal, y necesidad de mantenibilidad a largo plazo por parte de la empresa.

\section{Desarrollo por cortes verticales del sistema}
\label{sec:desarrollo-vertical-slices}

El desarrollo del aplicativo se organizó bajo el principio de cortes verticales o \textit{vertical slices}. Este enfoque evita construir primero una base de datos completa, luego una API completa y finalmente una interfaz completa sin conexión entre ellas. En su lugar, cada módulo se construye como una unidad funcional que atraviesa todas las capas: contrato de datos, validación, servicio de dominio, endpoint, cliente HTTP, hook de consulta, pantalla, permisos, pruebas y auditoría. Para el caso de CERMONT, este enfoque resulta adecuado porque cada paso del flujo de 14 etapas tiene reglas particulares, responsables y evidencias propias.

\begin{table}[!htbp]
\centering
\scriptsize
\caption{Estructura de un corte vertical aplicado a los módulos del aplicativo. Fuente: elaboración propia.}
\label{tab:vertical-slice-cermont}
\begin{adjustbox}{max width=\textwidth}
\begin{tabular}{|p{2.5cm}|p{4.2cm}|p{4.5cm}|p{4.0cm}|}
\hline\ClassicTableHeader
\textbf{Capa} & \textbf{Artefacto esperado} & \textbf{Propósito} & \textbf{Ejemplo en CERMONT}\\
\hline
Contrato compartido & Esquema y tipos de datos & Definir estructura de entrada y salida sin duplicidad. & Datos de orden, planeación o evidencia.\\
Backend & Ruta, controlador, servicio y modelo & Aplicar reglas de negocio y persistir cambios. & Crear paquete de planeación o registrar evidencia.\\
Frontend & Página, formulario, hook y estados UI & Permitir interacción del usuario con datos reales. & Vista de planeación o cierre administrativo.\\
Seguridad & RBAC, validación y auditoría & Restringir acciones y dejar rastro. & Solo ciertos roles aprueban planeación o cierre.\\
Pruebas & Unitarias, integración o funcionales & Verificar comportamiento esperado. & Rechazar ejecución si la planeación no está lista.\\
Documentación & Registro de decisión y evidencia & Explicar cómo se desarrolló el módulo. & Bitácora DOC asociada al módulo.\\
\hline
\end{tabular}
\end{adjustbox}
\end{table}

\section{Modelo de estados para controlar el avance operativo}
\label{sec:modelo-estados-control-avance}

El sistema no debe permitir que una orden avance por decisiones arbitrarias del usuario. En un proceso de cierre documentado, cada transición debe depender de precondiciones verificables. Por ejemplo, la ejecución no debería iniciar si la planeación no está aprobada; el informe no debería generarse si la ejecución no está terminada; la SES no debería radicarse si falta informe o acta; y la factura no debería emitirse si la SES no ha sido aprobada cuando el cliente la exige. Estas reglas transforman el aplicativo en un mecanismo de control operativo y no solo en una base de datos.

\begin{longtable}{|p{3.5cm}|p{4.2cm}|p{5.8cm}|}
\caption{Transiciones críticas propuestas para el flujo operativo-documental. Fuente: elaboración propia.}
\label{tab:transiciones-criticas-flujo}\\
\hline\ClassicTableHeader
\textbf{Transición} & \textbf{Precondición mínima} & \textbf{Razón de control}\\
\hline
\endfirsthead
\hline\ClassicTableHeader
\textbf{Transición} & \textbf{Precondición mínima} & \textbf{Razón de control}\\
\hline
\endhead
\hline
\multicolumn{3}{r}{Continúa en la siguiente página}\\
\endfoot
\hline
\endlastfoot
Solicitud a propuesta & Solicitud calificada y alcance preliminar definido. & Evita cotizar trabajos sin información mínima.\\
Propuesta a orden & Propuesta aprobada y PO registrada. & Garantiza que la ejecución tenga autorización comercial.\\
Orden a planeación & Orden activa y responsable asignado. & Permite preparar recursos, permisos y documentos.\\
Planeación a ejecución & Planeación aprobada y sin bloqueadores críticos. & Evita salida a campo sin herramientas, equipos o HES.\\
Ejecución a informe & Actividad finalizada y evidencias mínimas asociadas. & Evita informes incompletos o basados en memoria.\\
Informe a acta & Informe revisado y aprobado internamente. & Protege la entrega formal al cliente.\\
Acta a SES & Acta firmada o recibida por el cliente. & Soporta radicación administrativa.\\
SES a factura & SES aprobada por el cliente o plataforma correspondiente. & Evita facturación sin soporte aprobado.\\
Factura a pago & Factura aprobada y referencia de pago registrada. & Permite cierre administrativo trazable.\\
\end{longtable}

\section{Diseño del módulo de archivos, evidencias y firmas}
\label{sec:diseno-modulo-archivos-evidencias}

El módulo de evidencias tiene una responsabilidad crítica: transformar fotos y documentos dispersos en soportes consultables, asociados y auditables. En la práctica operativa, una evidencia sin contexto pierde valor: una fotografía enviada por mensajería instantánea puede mostrar una actividad, pero si no se sabe a qué orden pertenece, quién la tomó, cuándo se capturó y qué componente documenta, su utilidad para el cierre técnico y administrativo disminuye. Por ello, el sistema debe almacenar metadatos junto con cada archivo.

Los metadatos mínimos incluyen identificador de orden, tipo de evidencia, usuario responsable, fecha, estado de sincronización, nombre seguro de archivo, tipo MIME, tamaño, observaciones y relación con el paso del proceso. En caso de firmas, se debe conservar vínculo con el firmante, rol, fecha y documento asociado. Esta estructura responde a la falla documentada de evidencias dispersas y fortalece la generación posterior de informes y actas.

\section{Diseño del módulo de costos reales}
\label{sec:diseno-modulo-costos-reales}

El control de costos reales se concibe como una capacidad transversal y no como una pantalla aislada de contabilidad. La propuesta económica define una estimación; la planeación anticipa recursos; la ejecución registra consumos; el cierre administrativo consolida factura y pago. Si estos datos permanecen desconectados, la empresa no puede analizar desviaciones. El diseño del módulo de costos debe permitir que cada orden acumule líneas de costo por materiales, mano de obra, herramientas, equipos, transporte, subcontratos, impuestos configurables e imprevistos.

Esta funcionalidad debe implementarse con prudencia académica. Si no existen datos reales cargados, el libro debe describir el diseño y los criterios de medición, pero no reportar márgenes, variaciones o ahorros. La contribución del proyecto en esta fase consiste en dejar preparada la estructura para comparar costos estimados y reales cuando CERMONT alimente el sistema con información operativa verificable.

\section{Integración entre backend, frontend y documentación técnica}
\label{sec:integracion-backend-frontend-docs}

La integración entre capas se construye alrededor de la coherencia de nombres, contratos y responsabilidades. El backend valida, aplica reglas de negocio y persiste. El frontend presenta, captura datos y comunica estados. Los paquetes compartidos definen contratos y tipos comunes. La documentación técnica actúa como referencia para evitar duplicidad de lógica, rutas huérfanas, pantallas sin endpoints o módulos que solo muestran títulos. Esta disciplina responde a una observación recurrente en auditorías de software: una aplicación puede compilar y aun así no resolver el proceso de negocio si sus módulos no están conectados con datos reales y reglas verificables.

Por esta razón, el libro debe explicar no solo qué tecnologías fueron usadas, sino cómo interactúan en cada módulo. En una orden de trabajo, por ejemplo, el frontend no debe inventar estados de avance; debe consultar una proyección o entidad calculada desde backend. El backend no debe aceptar cualquier transición; debe verificar precondiciones. Los contratos no deben duplicarse manualmente; deben mantenerse como fuente común. Esta articulación convierte la arquitectura en respuesta directa a la problemática y no en un listado de herramientas.
